\section{Methods}\label{sec:methods}

All four analyses were metadata- and documentation-only. No measured metabolite value, data matrix, raw or derived
data file was opened at any stage. Each analysis was specified in a written protocol, implemented in code with an
offline test suite, and bound by SHA-256 hashes in a lock file signed by the author before execution. The locks were
not registered with an external registry. Protocols, code, locks and outputs are public
(Declarations).

\subsection{Target property and evidence standard}\label{sec:target}

The unit of interest is a deposited assay: the biological samples measured on one platform for one matrix. We ask
whether the deposited metadata or linked documentation identify its independent sampling units: whether the samples
come from human participants, how many distinct participants contributed, and whether any contributed more than one
sample (time points, pre/post samples, several specimens, replicates entered as biological samples) or, failing that,
whether a sample-to-participant mapping lets an analyst model the dependence. Evidence had to be affirmative: unique
sample or source labels were bookkeeping, not evidence; counts had to be stated integers or sums over groups the source
called disjoint; an enrolment figure counted only if the text tied it to the analysed samples. Independence established
in this sense is ``documented''; a study whose documentation is silent may still consist of independent participants.

\subsection{MetaboLights census}\label{sec:census}

A locked enumeration audit fixed the universe at 3,582 accessions (REST list $\cap$ FTP mirror, minus two
development exemplars) on 25 October 2026. The census read only investigation, sample and assay files; data files were
refused by an allow-list guard. Studies were screened in a fixed order (human organism; removal of quality-control,
blank, pooled and reference rows; sample-name integrity; a single matrix; the independence test). The independence test
had two affirmative routes: a person-identifier column with exactly one sample per identifier (R1), or a
cross-sectional design annotation plus a study-description sentence matching one of four frozen
one-sample-per-participant templates, guarded against negation (R2); both required at least 50 sources. The strict
criterion (L1) counted R1 and explicit R2 statements; a sensitivity count added count-only statements. Studies that
failed the test but had at least 50 distinct source labels with one sample each, or at least 50 samples with a
degenerate source column, became unresolved candidates for documentary assessment (Supplementary Methods S1).

\subsection{Secondary audit of Metabolomics Workbench subject fields}\label{sec:workbench}

This analysis is secondary. It reuses 104 human Metabolomics Workbench studies retained on 17 October 2026 by an
earlier value-blind screen built for a different, subsequently closed research question; its criteria and exclusion
counts, copied from its hashed closure record, are in Supplementary Table S1. The 104 are not a representative sample
of the Workbench, and no inference is drawn about the repository as a whole. For each study we read the
\texttt{factors} table returned by the Workbench REST interface (from the screen's cache; no network request) and
recorded whether any factor key named a subject, patient, participant, individual, donor, person, volunteer or proband.
Workbench mwTab files, which may hold an optional subject column, embed measured values and were not opened, so nothing
is claimed about subject fields elsewhere in the Workbench.

\subsection{Documentary assessment of a random tranche (T1)}\label{sec:t1}

\paragraph{Population and selection.} The census produced 371 unresolved MetaboLights candidates with at least 50
unique source labels (312) or at least 50 samples with a degenerate source column (59). A tranche of 40 was selected
as follows. Each candidate was ranked by SHA-256 of a salt concatenated with its accession, and the first 40 were
taken. The salt was the output value of a NIST Randomness Beacon (version 2) pulse, at a time fixed in the signed lock
before the pulse existed (pulse 2-2005958, 30 October 2026, 15:00 UTC). No candidate was replaced. The tranche is
therefore a simple random sample of 40 of the 371 candidates, and finite-population inference applies.

\paragraph{Sources and document freeze.} Extractors used deposited investigation files and restricted reading views of
linked or accession-matched publications. The views excluded results-type material. No author was contacted. The
author signed an inventory of document and view hashes before extraction, after which no source could be added
(Supplementary Methods S3).

\paragraph{Questions and grades.} For the selected assay, extractors recorded human origin (Q1), assay membership
(Q2), participant count and scope (Q3), one-sample or repeated sampling (Q4), overlap flags (Q6), and any published
sample-to-participant mapping (Q7). Each answer other than \emph{not stated} required a located verbatim quotation. Deterministic
grading followed Supplementary Table S2. STRICT requires one sample per participant for a downstream analysis;
Q4 and Q7 separately address sampling structure and reconstructible linkage.

\paragraph{Extraction, grading and adjudication.} Extraction was performed by AI sub-agents working under a written
brief fixed and hashed before any record was produced (Section~\ref{sec:ai}). The extractors were not told the
deposits' biological row counts; the correspondence between a stated count and the deposited rows (rules 6--7) was
computed only by the grading function. Every quotation was checked automatically for verbatim presence in the cited
view. The author adjudicated every CONTRADICTED and STRICT record against the filed evidence; a rejected
CONTRADICTED record became UNRESOLVED. No record was edited after grading.

\subsection{Pre-verification of deposits removed by quality-control filtering}\label{sec:a35}

A diagnostic after the census identified 19 MetaboLights deposits with human organism values and at least 50 sample
rows in which every row had been removed by the M2 quality-control rule. Because the census selected no assay for
them, the four checks were applied to the biological samples listed in each deposit's sample file: V1, human
participants; V2, the rows are identified as participants' samples by fields other than the sample-type label, or by
a publication; V3, an integer number of biological samples stated in text (counting rows was not accepted); and V4,
the sampling structure, failing on explicit repeated sampling. A study was verified if all four checks passed and
failed if any check failed. Sources were the investigation, sample and assay files and named publications. The
evidence set was frozen and signed before reading, extraction was blind to row counts, and quotations of sample
names or individual-level rows were prohibited.

\subsection{Statistical analysis}\label{sec:stats}

For the random tranche, the number of STRICT studies among $n = 40$ drawn without replacement from $N = 371$ bounds
the number $K$ of documentation-recoverable candidates: the one-sided 95\% upper bound is the largest $K$ with
$P(X \le x \mid N, K, n) > 0.05$ under the hypergeometric distribution (Clopper--Pearson bounds scaled by $N$ are given
for comparison). These bound candidates whose documentation meets the standard, not truly independent cohorts. All
other results are counts with their denominators.

\subsection{Author review of the AI-assisted extraction}\label{sec:validation}

The review protocol was hashed on 31 October 2026, before selecting studies. The intended independent reviewer
could not be recruited, so the author performed both components on 1--2 November 2026. In component A the author
re-extracted Q3 and Q4 for 10 of the 17 studies with publication text that the author had not adjudicated, selected by
prespecified hash ranking. The author had not seen the original answers or quotations but knew the classifications
and had seen a second model's answers for the same studies. Agreement is reported by field and is not extrapolated.
In component B the original adjudicator checked all 12 CONTRADICTED records for explicit supporting statements.
The review did not change the frozen grades. Supplementary Methods S5 records the deviations, including the lack of
per-study completion times.

\subsection{Use of AI tools}\label{sec:ai}

Evidence extraction for the tranche and the pre-verification was performed by sub-agents of Claude (Anthropic),
working under fixed, hashed written briefs. Each sub-agent read only the sealed files of its assigned studies,
produced structured records with verbatim quotations, and had to pass an automated validator that checked the schema,
the verbatim presence of every quotation and the form of every count. Grades were assigned by deterministic code, not
by the sub-agents. The author adjudicated every CONTRADICTED record and carried out the human validation in
Section~\ref{sec:validation}. The briefs, validators and raw sub-agent outputs are public. Claude also assisted with
code and drafting under the author's direction; the author is responsible for the content.

